\begin{hbconvention}{def:p06-def-empty-base}{Constructibility at the empty base}
\emph{Editorial correction to the preceding exact source.}
The handbook's satisfaction relation is defined only for structures with
nonempty domains.  Accordingly, read the preceding definition with the
eliminable base clause
\[
  \operatorname{Def}(\varnothing)\defeq\{\varnothing\},
\]
and use its satisfaction-based display only when \(A\ne\varnothing\).  Thus
\(L_0=\varnothing\), \(L_1=\{\varnothing\}\), and
\(L_{\alpha+1}=\operatorname{Def}(L_\alpha)\) uses the displayed semantic
clause from the first nonempty stage onward.  This convention is necessary to
type \(L_1\); it stands as a base case for the existing recursion, differs
from permitting empty first-order structures globally, and is used in the
later relativization proof.  There, the sentence saying that Separation works
“whatever the satisfaction clause returns” is to be read instead as the
direct calculation \(\varnothing\in L_1\subseteq L\).
\end{hbconvention}
